\pdfbookmark{Общая характеристика работы}{characteristic}             % Закладка pdf
\section*{Общая характеристика работы}

\newcommand{\actuality}{\pdfbookmark[1]{Актуальность}{actuality}\underline{\textbf{\actualityTXT}}}
\newcommand{\progress}{\pdfbookmark[1]{Разработанность темы}{progress}\underline{\textbf{\progressTXT}}}
\newcommand{\aim}{\pdfbookmark[1]{Цели}{aim}\underline{{\textbf\aimTXT}}}
\newcommand{\tasks}{\pdfbookmark[1]{Задачи}{tasks}\underline{\textbf{\tasksTXT}}}
\newcommand{\aimtasks}{\pdfbookmark[1]{Цели и задачи}{aimtasks}\aimtasksTXT}
\newcommand{\novelty}{\pdfbookmark[1]{Научная новизна}{novelty}\underline{\textbf{\noveltyTXT}}}
\newcommand{\influence}{\pdfbookmark[1]{Практическая значимость}{influence}\underline{\textbf{\influenceTXT}}}
\newcommand{\methods}{\pdfbookmark[1]{Методология и методы исследования}{methods}\underline{\textbf{\methodsTXT}}}
\newcommand{\defpositions}{\pdfbookmark[1]{Положения, выносимые на защиту}{defpositions}\underline{\textbf{\defpositionsTXT}}}
\newcommand{\reliability}{\pdfbookmark[1]{Достоверность}{reliability}\underline{\textbf{\reliabilityTXT}}}
\newcommand{\probation}{\pdfbookmark[1]{Апробация}{probation}\underline{\textbf{\probationTXT}}}
\newcommand{\contribution}{\pdfbookmark[1]{Личный вклад}{contribution}\underline{\textbf{\contributionTXT}}}
\newcommand{\publications}{\pdfbookmark[1]{Публикации}{publications}\underline{\textbf{\publicationsTXT}}}
%%%%%%%%%%%%%%%%%%%%
\newcommand{\KKT}{\pdfbookmark[1]{Курило-Камчатский желоб}{kkt}\underline{\textbf{\kktTXT}}}
\newcommand{\AT}{\pdfbookmark[1]{Алеутский желоб}{at}\underline{\textbf{\atTXT}}}
\newcommand{\MAT}{\pdfbookmark[1]{Центральноамериканский желоб}{mat}\underline{\textbf{\matTXT}}}
\newcommand{\PCT}{\pdfbookmark[1]{Перу-Чилийский желоб}{pct}\underline{\textbf{\pctTXT}}}
\newcommand{\KT}{\pdfbookmark[1]{Кермадек желоб}{kt}\underline{\textbf{\ktTXT}}}
\newcommand{\TT}{\pdfbookmark[1]{ЖелобТонга желоб}{tt}\underline{\textbf{\ttTXT}}}
\newcommand{\VT}{\pdfbookmark[1]{Вануату желоб}{vt}\underline{\textbf{\vtTXT}}}
\newcommand{\VTT}{\pdfbookmark[1]{Витязь желоб}{vtt}\underline{\textbf{\vttTXT}}}
\newcommand{\NBT}{\pdfbookmark[1]{Новобританский желоб}{nbt}\underline{\textbf{\nbtTXT}}}
\newcommand{\SCT}{\pdfbookmark[1]{Сан-Кристобаль желоб}{sct}\underline{\textbf{\sctTXT}}}
\newcommand{\PT}{\pdfbookmark[1]{Филиппинский желоб}{pt}\underline{\textbf{\ptTXT}}}
\newcommand{\MNT}{\pdfbookmark[1]{Манильский желоб}{mnt}\underline{\textbf{\mntTXT}}}
\newcommand{\PLT}{\pdfbookmark[1]{Палау желоб}{plt}\underline{\textbf{\pltTXT}}}
\newcommand{\YT}{\pdfbookmark[1]{Яп желоб}{yt}\underline{\textbf{\ytTXT}}}
\newcommand{\RT}{\pdfbookmark[1]{Рюкю желоб}{rt}\underline{\textbf{\rtTXT}}}
\newcommand{\MT}{\pdfbookmark[1]{Марианский желоб}{mt}\underline{\textbf{\mtTXT}}}
\newcommand{\IBT}{\pdfbookmark[1]{Идзу-Бонин желоб}{ibt}\underline{\textbf{\ibtTXT}}}
\newcommand{\JT}{\pdfbookmark[1]{Японский желоб}{jt}\underline{\textbf{\jtTXT}}}
\newcommand{\PST}{\pdfbookmark[1]{Пюйсегюр желоб}{pst}\underline{\textbf{\pstTXT}}}
\newcommand{\HT}{\pdfbookmark[1]{Хикуранги желоб}{ht}\underline{\textbf{\htTXT}}}


{\actuality} Актуальность диссертации заключается в ее междисциплинарном характере, демонстрирующем сочетание из следующих подходов: 

\begin{enumerate}[label=\fbox{\arabic*}]
	\item систематический анализ геопространственных данных из множества источников; 
	\item статистическое моделирование и обработка данных библиотеками Python и R, AWK и Octave; 
	\item геологический анализ литературы; 
	\item картографическое отображение и моделирование желобов скриптами Generic Mapping Tools (GMT) и плагинами QGIS. 
\end{enumerate}
Диссертация вносит вклад в исследования по геотектонике, геологии и геоморфологии дна Тихого океана. 

Формирование глубоководного желоба представляет собой сложный процесс, на который влияет совокупность различных геологических и тектонических факторов, объясняющих процесс его формирования и дальнейшее развитие в зависимости от степени влияния данных факторов, влияющих на морфоструктуру. Попытки исследований желобов по поперечному сечению профилей проводились в прошлом, например, но все же основанные на технических приемах того времени. На настоящий момент в существующей литературе нет работ, основанных на систематическом применения скриптов GMT для моделирования морфоструктуры океанских желобов, и это исследование способствовало этому вопросу.

В настоящее время изучение морских геологических явлений и сложных процессов с помощью программирования и скриптов является эффективным методом. Быстрое развитие передовых методов анализа данных представляет такие инструменты, как GMT, Octave/MATLAB, R и Python. Использование скриптовых методов особенно эффективно применительно к огромному количеству морских геологических данных (т.н. big data). Обработка больших массивов данных с помощью скриптовой техники картографирования и графической визуализации представляется критически важным подходом, поскольку алгоритмы библиотек предоставляют доступ к точному и быстрому анализу больших массивов данных. Конкретная информация о глубинах труднодоступных океанических желобов может быть получена для точной визуализации и анализа их подводной геоморфологии от локального до регионального и глобального масштабов.

Однако, несмотря на существующие отдельные публикации, в настоящий момент отсутствует единообразие в изучении глубоководных желобов, образовавшихся в определенных геологических и тектонических условиях на окраинах Тихого океана, а также имеет место нехватка систематического картографирования желобов применительно к Тихоокеанскому региону. Отсутствие системного подхода к картографированию глубоководных желобов приводит к недостаточному пониманию геоморфологических различий их профилей в разных частях океана: южном и северном, восточном и западном, а также и их особенностей, вызванных геологическими и тектоническими локальные условиями в местах формирования. 

В связи с этим, данная диссертация направлена на восполнение существующих на данный момент методологических пробелов путем систематического картографирования, геоморфологическое моделирования, пространственного анализа и классификации 20 желобов Тихого океана. В частности, данная работа представляет собой междисциплинарный подход, применяя сочетание следующих дисциплин для решения отдельных задач: 

\begin{enumerate}[label=(\roman*)]
	\item геодинамика и тектоника (анализ происхождения, эволюции и формирования желобов); 
	\item картография (скриптовые алгоритмы обработки пространственных данных в GMT и плагины QGIS); 
	\item статиcтическая обработка данных (библиотеки Python, R и Octave); 
	\item геоморфология (анализ современных морфоструктур рельефа дна Тихого океана). 
\end{enumerate}

В диссертации представлен системный подход к картографированию, моделированию и сравнительному анализу морфоструктур глубоководных желобов на основе сочетания передовых технических методов визуализации и анализа данных в области картографии и статистики и теоретических знаний по геодинамике и эволюции глубоководных желобов Тихого океана.

% {\progress}
% Этот раздел должен быть отдельным структурным элементом по
% ГОСТ, но он, как правило, включается в описание актуальности
% темы. Нужен он отдельным структурынм элемементом или нет ---
% смотрите другие диссертации вашего совета, скорее всего не нужен.

{\aim} данной работы является сравнительный анализ вариативности геоморфологии 20 глубоководных желобов Тихого океана в их различных сегментах на основе прикладного статистического анализа, поддерживаемого GMT, Octave, Библиотеки Python и R, и геоморфологическое моделирование желобов для классификации и ранжирования их форм рельефа на основе общих характеристик в ортогональном плане:
\begin{enumerate}
    \item Проанализировать вариации геоморфологической формы подводной геоморфологии в 20 глубоководных районах желоба Тихого океана. 
    \item Выполнить классификацию геоморфологических типов 20 глубоководных желобов на основе их геометрических свойства, отличия и сходство. 
\end{enumerate}
Следующие типы форм геоморфологических профилей желобов были выявлены и подразделены на семь типов: U-образные (в плане), V-образные (в плане), асимметричный, серповидный, извилистый, удлиненный, каскадный. Для каждого типа (U, V, асимметричный и т. д.) выделены характерные подтипы крутизны: сильная, очень сильная, экстремальная, крутая, очень крутой. Склоны долин классифицируются следующим образом: очень высокие, высокие, умеренные, низкие, в зависимости от кривизны и степени. Классы размеров и уклонов долин анализируются в контексте физической среды, тектонических и геологических условий района желоба.

Для~достижения поставленной цели необходимо было решить следующие {\tasks}:
\begin{enumerate}[label=\Roman*.]% https://tex.stackexchange.com/a/476052/104425
  \item Моделирование морфологии продольно-поперечных разрезов (профилей) глубоководных желобов с помощью модулей скриптового картографического инструмента GMT для систематического построения ряда профилей поперечного сечения для сравнительного анализа геоморфологии 20 желобов.
  \item Картографирование локальных топографических и геологических условий в районах расположения 20 глубоководных желобов, на основе имеющихся наборов данных.
  \item Картографирование геофизических свойств (морской гравитации в открытом воздухе и модели геоида) 20 желобов на основе имеющихся наборов данных. 
  \item Выполнение статистического анализа батиметрических продольно-поперечных разрезов (профилей) глубоководных желобов на основе использования методов алгоритмов программирования (Python и R) для получения моделей взаимосвязи геологических факторов и геоморфологии желобов.
  \item Методическая разработка совместного применение статистических библиотек для анализа и визуализации данных (Python, R, GMT, Octave, AWK, QGIS) для анализа данных и моделирования
  \item Анализ вариативности морфологии желобов вдоль окраин Тихого океана для оценки взаимосвязи геодинамических факторов, влияющих на их морфоструктуру. 
\end{enumerate}

{\novelty}
\begin{enumerate}[beginpenalty=10000] % https://tex.stackexchange.com/a/476052/104425
  \item Впервые разработаны экспериментальные методы GMT, примененные к моделированию именно глубоководных желобов, а также расширении существующих методов пространственного моделирования с использованием языков программирования Python, R, AWK и Octave для моделирования, статистического анализа, визуализации и геоморфологической классификации подводных форм рельефа желобов. Использование методов расширенного анализа данных оказало решающее значение для полученной точной и надежной обработки данных, поскольку понимание форм рельефа морского дна может быть основано исключительно на компьютерном моделировании данных из-за их недоступности месторасположения желобов. 

  \item Техническая инновация работы заключается в междисциплинарном подходе, сочетающем ГИС-анализ и статистические методы (R, Python, MATLAB), которые способствуют изучению океанских желобов, при пространственном анализе больших данных. Моделирование закономерностей корреляции геоморфологической структуры и геологических особенностей морского дна в районе желобов имеют решающее значение для моделирования рельефа дна океана. Выбор методологии, инструментов и алгоритмов объясняется целями и задачами исследования, поскольку специфика морской геологии, геодинамики и геофизики заключается в особо высоких требованиях к методическому уровню обработки данных, применении самых передовых технологий и инструментария для точного анализа данных, недоступных для непосредственного изучения в связи с удаленностью объекта (глубоководные желоба). 
  
\item Было выполнено оригинальное исследование на основе наборов данных, которые обрабатывались, вычислялись и анализировались в полуавтоматическом режиме с помощью алгоритмов машинного обучения в картографии. Таким образом, диссертация представляет комплексный подход. Основанные на глубинном анализе геологического строения желобов (представленный в описательной главе 2) и технические подходы, использующие современные алгоритмы анализа данных и эффективную визуализацию путем использования языков программирования и скриптовой картографической программы GMT, не применяемого ранее в существующих работах по сравнительной морфологии глубоководных желобов Тихого океана. 

\item Впервые представлена систематическая классификации и осуществлено сравнительное моделировании геоморфологических профилей глубоководных желобов Тихого океана, выполненных с помощью последовательного использования современных наборов инструментов для создания картографических скриптов и статистического анализа. Техническая новизна заключается в сочетании GIS, GMT, Python, AWK, R применительно к задачам картографирования и моделирования подводных форм рельефа дна глубоководных желобов Тихого океана.  
\end{enumerate}

{\influence} Геоморфология рельефа дна Мирового океана в целом и океанических желобов в частности (как особых его структур), их образование, формирование и эволюция – вопрос особой важности для научного сообщества в области геодинамики, тектоники, морской геофизике и геологии. Информация о геоморфологии глубоководных  желобов имеет важное значение для изучения динамической обстановки и тектонической активности окраин Тихого океана. Актуальность этого вопроса значительно увеличилось с начала бурного развития информационных технологий (ИТ), инструментов и методов расширенного анализа данных, но их понимание тем не менее остается неоднородным. Поскольку большинство желобов расположены по окраинам именно Тихого океана (по сравнению с другими), изучение тихоокеанских желобов играет центральную роль для анализа и понимания их образования. 

Особые геологические условия, наличие зон тектонической субдукции, обширная территория Тихого океана со сложной системой циркуляции (Рис. 1), распространение зоны «огненного кольца», т.е. сейсмически активного пояса землетрясений и вулканических процессов, делают желоба Тихого океана чувствительными к факторам, влияющим на их формирование, которое вызывает изменения в их геоморфологических формах. До 80\% землетрясений, зафиксированных сейсмическими службами, сосредоточены в Тихоокеанском «огненном кольце» [6, 7]. Тектонические структуры дня Тихого океана отличаются вариативностью и сложностью форм. В общей классификации можно выделить острова, подводные окраины материков, океанические валы, подводные горы, срединно-океанические хребты, котловины, каньоны, абиссальные равнины, ложе океанических впадин, подводные внутриокеанические каньоны, а также глубоководные желоба, зоны разломов, срединноокеанические хребты и срединноокеанические рифты, которые входят в единую систему рифтов растяжения. В этом контексте, наиболее репрезентативными индикаторами вариаций глубоководных желобов являются геологические и тектонические факторы, такие как динамика литосферной коры, влияющая на скорость и интенсивность субдукции плит, величину и повторяемость подводных вулканов, вызывающих активную седиментацию. 

Геоморфология морского дна в желобах зависит от множества факторов, которые в конечном счете влияют на форму их рельефа, определяют вариативность морфоструктур и динамику. Анализ геоморфологии желобов имеет практическое значение при изучении тектонически активных зон Тихого океана. Эта диссертация рассматривает такие факторы как формирование тектонических плит, субдукцию, историческое геологическое развитие, расположение основных зон землетрясений и подводных вулканов как основные типы ключевых факторов, влияющих на формирование океанических желобов. Вторичные факторы включают океанические течения, седиментацию и особенности распределение биоты, способствующие седиментации. 

Знание и правильное понимание движущих факторов, воздействующих на геоморфологию океанического дна, дает понимание геоморфологии океанского дна. Учитывая важность глубоководных районов, изучение морфоструктур желобов является актуальной задачей. Используя методы моделирования данных, были сопоставлены и проанализированы формы пересечения профилей 20 желобов в ортогональном направлении для выделения различий и вариации в их геоморфологии. 

{\methods} Мультидисциплинарная интеграция картографических методов исследования (GIS, GMT), статистической обработки данных (Python, AWK, R) и анализа обширного материала литературных источников по геологии и тектонике Тихого Океана применительно к задачам картографирования и моделирования подводных форм рельефа дна глубоководных желобов Тихого океана.

{\defpositions}
\begin{enumerate}[label=\emph{\alph*})]
 	 \item Первое положение. Научное обоснование, комплекс картографических методических решений и скриптовых алгоритмов GMT для моделирования морфоструктур глубоководных желобов Тихого океана, находящихся под воздействием региональных тектонических условий и меняющихся скоростей движения литосферных плит.
  	\item Второе положение. Основные пути повышения: i) методической актуальности (численные методы, скриптовые языки и программирование на языках Python, R, Octave); ii)  наукоемкости (анализ обширного объема материалов как картографического, так и текстуального типа); iii) информативности (обработка больших массивов геолого-геофизических и геоморфологических данных с неочевидными, нелинейными связями, big data) для комплексных исследований дна океана при совокупном влиянии тектонических и геодинамических факторов, определяющих морфоструктуру дна Тихого океана.
 	 \item Третье положение. Теоретическая аргументация методологических решений по созданию методики комплексного картографирования профилей желобов двух- и трехмерным моделированием (2D и 3D-моделирование), описывающих: а) изменение крутизны и уклона профилей при вариативности геологических и тектонических условий; б) изменение геофизических свойств области исследования при вариациях гравиметрии вдоль русла желобов.
 	 \item Четвертое положение. Современный геоинформационных методический комплекс с применением скриптовых решений GMT и статистического анализа на языках Python, R и Octave для исследований изменения морфологических, геометрических и геофизических свойств желобов под воздействием геолого-геофизических региональных эффектов от воздействии вариативности скорости движения и субдукции литосферных плит.
  	 \item Пятое положение. Западный и восточный регионы Тихого океана характеризуются вариативностью типов глубоководных желобов и маркированы U, V-образными и асимметричными формами разрезов поперечного сечения. Сейсмически активные зоны и геолого-геофизические условия определяют морфологию желобов.
\end{enumerate}

%{\reliability} полученных результатов обеспечивается \ldots \ Результаты находятся в соответствии с результатами, полученными другими авторами.

{\probation}
По результатам исследований, выполненных в рамках данной диссертации, опубликовано 70 основных печатных работ, из них 21 статей в периодических изданиях, индексированных в Scopus и Web of Science и входящих в список рекомендованных ВАК, 3 тезиса докладов материалов конференций и 46 публикаций в прочих журналах.

{\contribution} Автор принимала активное непосредственное участие в работе.
\begin{enumerate}[label=\emph{\Alph*})]
	\item Автор лично написала скрипты GMT, использованные для моделирования желобов и тематического картографирования, а также скрипты на языках программирования Python, R и Octave для дополнительного статистического анализа батиметрии желобов. С использованием данных скриптов автором лично проведено картографирование, 2D и 3D моделирование всех основных 20 желобов, представленных в настоящей диссертации, геологическое описание исследованных районов желобов Тихого океана. 

	\item Автором лично выполнены все этапы картографирования, включающие следующие шаги: сбор, обработка, анализ, интерпретация и визуализация данных и моделирования желобов. Автором лично выполнен отбор и предварительный анализ данных на пригодность (разрешение, надежность источника, актуальность) из доступных геофизических, геологических и топографических источников данных. Автором лично проведена обработка данных для картографирования (форматирование, проецирование и подготовка отдельных слоев по координатам профилей) и последующее систематическое картографирование. 

	\item Все иллюстрации без исключения (как карты, так и графики), представленные в диссертации, являются результатом личной работы автора: картографирование на основе GMT и частично QGIS с использованием плагинов для цифрования, графики профилей желобов, оцифрованные в GMT по авторской разработанной методике с использованием модулей GMT, статистические графики, выполненные на основе R, Python и Octave. При осуществлении геоморфологического моделирования использовалась авторская разработка профилирования желобов с помощью комбинации модулей GMT, позволяющая с высокой скоростью и большой точностью производить оцифровку профилей по параметрам, разработанным автором, представленным в данной диссертации и апробированным на 20 желобах Тихого океана.

	\item По выполнении картографирования и технических скриптовых работ, автор лично провела интерпретацию морфологии желобов на основе результатов моделирования, а также анализ вариативности геофизических и геологических условий в областях исследования (мест расположения желобов) на основе визуализации данных. В соответствии с целью и задачами диссертации, автор проанализировала и обобщила обширный фактический материал (как картографические, так и литературные источники), опубликованные в предыдущих существующих работах. Результаты графического моделирования и картографирования были использованы автором как иллюстрации в публикациях. 
\end{enumerate}

\ifnumequal{\value{bibliosel}}{0}
{%%% Встроенная реализация с загрузкой файла через движок bibtex8. (При желании, внутри можно использовать обычные ссылки, наподобие `\cite{vakbib1,vakbib2}`).
    {\publications} Основные результаты по теме диссертации изложены
    в 70 печатных изданиях, 21 из которых изданы в журналах, рекомендованных ВАК, Scopus и Web of Science,
    3 тезиса докладов материалов конференций и 46 публикаций в прочих журналах.
}%

 % Характеристика работы по структуре во введении и в автореферате не отличается (ГОСТ Р 7.0.11, пункты 5.3.1 и 9.2.1), потому её загружаем из одного и того же внешнего файла, предварительно задав форму выделения некоторым параметрам

%\underline{\textbf{Финансирование}}. Диссертационная работа была выполнена при поддержке гранта No. 0144-2019-0011 ИФЗ им. О. Шмидта РАН.

\underline{\textbf{Объем и структура работы.}} Диссертация состоит из~введения,
5 глав, заключения и~приложения. Полный объем диссертации
\textbf{269}~страниц текста с~\textbf{87}~рисунками и~7~таблицами. Список
литературы содержит \textbf{351}~наименование.

\pdfbookmark{Содержание работы}{description}                          % Закладка pdf
\section*{Содержание работы}
Во \underline{\textbf{введении}} обосновывается актуальность исследований как для геологической картографии, так и для теоретических изысканий в области тектоники и геоморфологии дня, проводимых в~рамках данной диссертационной работы, далее 
приводится обзор научной литературы по~изучаемой проблеме (включая как российских, так и зарубежных авторов), сформулирована цель, поставлены задачи работы, излагается научная новизна и практическая значимость представляемой работы для геолого-геоморфологических исследований Тихого Океана (Рис. \ref{fig:01}). В~последующих главах описываются общие принципы методологии скриптовой картографии, использованные в данной работе, затем представлена апробация на частных примерах 20 глубоководных желобов Тихого Океана.

\underline{\textbf{Первая глава}} посвящена обзору литературы и разработанности темы.

История изучения геоморфологии и геологии дна Мирового океана постоянно развивается на протяжении 20-го и 21 веков, что отражено в многочисленных публикациях. Вопрос образования и строения океанических желобов издавна привлекал внимание. Исследования желобов проводились в различных аспектах: батиметрическая съемка для оценки объемов минеральных ресурсов, прогнозирование и частотность землетрясений. Понимание разнообразия форм рельефа дна – сложная и наукоемкая задача, основанная на сочетании междисциплинарных подходов: пространственная визуализация данных, статистический анализ, геологическое моделирование, применение численных методов и пакетной обработки массивных данных. 

ХХ-й век отмечен систематическим исследованием океана с использованием методов ДЗЗ, ГИС для обработки данных, инструментов машинного обучения для статистической обработки и анализа данных, геолого-геофизическим исследований. На сегодняшний день существуют и постоянно развиваются разные подходы к картографированию и мониторингу морского дна. Картографирование особенностей подводных морфоструктур и понимание факторов их вариативности объединяет несколько важных компонентов в развитии морской геологии: 1) поддержка пространственных морских наблюдений и управления массивными наборами данных; 2) мониторинг Мирового океана; 3) получение знаний о форме и процессах Земли; 4) оценка запасов минеральных ресурсов в глубинах океана; 5) обновление знаний и данных о геотектонике и геодинамике применительно к областям океанского дна.

Стремительный технический прогресс в XXI в. отмечен развитием ИТ-технологий, что позволило более эффективно анализировать геологические и океанологических данные используя концептуальный подход “больших данных” (big data), что требует значительной компьютеризации и машинных методов. Это существенно улучшило качество и скорость обработки данных в геологической картографии. В XXI в. к преимуществам применения современных ИТ-технологий для геотектоники и геодинамики относится возможность применять языки программирования для моделирования процессов динамики Земли и анализа данных, а также открытые данные). XXI в. открыл новые возможности для междисциплинарных подходов в геонауках, которые включают в себя скриптовые алгоритмы обработки данных в картографии, пространственный анализ, ГИС, машинная обработка и визуализация данных, статистический анализ атрибутивных данных на основе библиотек языков программирования высокого уровня, таких как Python и R. 

Начиная с 2000 г, в некоторых работах представлено использование Python и R, скриптовое картографирование и машинные методы анализа данных. Тестируя функциональность пакетов, алгоритмов и методов, эти работы бросили вызов науке о данных и представили новые методы исследований в геологии. Однако большая часть существующих геологических исследований по океаническим желобам сосредоточены на проблеме субдукции тектонических плит и движениях литосферы в основном применяя традиционные методы ГИС. 

Несмотря на очевидную важность междисциплинарной связи между теоретическими изысканиями в области тектоники и геодинамике и практическими применениями скриптовых алгоритмов обработки данных, на данный момент полностью отсутствуют исследования, объединяющие эти области в едином подходе. Данная диссертация восполняет этот пробел путем применения практических методов высокоуровневой скриптовой картографии для теоретических изысканий в области геодинамики и тектоники на примере Тихого океана. 

Моделирование глубоководных желобов всегда ограничивалось как их удаленностью, так и сложностью доступа к высококачественным данным по крупномасштабной съемке морского дна. Тем не менее, в наблюдениях Земли методами данных ДЗЗ происходит непрерывный технический прогресс и данные улучшаются в качестве и доступности. Совершенствование методов обработки данных и развитие скриптовых языков позволяют точно отображать и моделировать наименее доступные объекты дна в самых удаленных частях желобов. Использование этих подходов для анализа морфоструктур желобов позволяет оценить вариативность их форм, классифицировать и выделить элементы подводного рельефа. Скриптовые методы произвели революцию в геологическом картографировании, позволив собирать контекстную информацию для визуализации дна океана. 

Геоморфология океанических желобов картографирована существенно хуже и менее детально по сравнению с наземными участками Земли. Это связано прежде всего с техническими трудностями регулярной съемки морского дна. На основании выполненного обзора литературы автором выявлено, что несмотря на существующие разрозненные работы по отдельным желобам, системное картографирование морфоструктур желобов Тихого океана в литературе на данный момент отсутствует, что требует отдельной объемной работы. Представленная диссертация представляет собой восполнение это проблемы. 

\begin{figure}[ht]
    \centerfloat{
            \includegraphics[width=0.85\linewidth]{Рис_01.jpg}
    }
    \caption{Топографическая (батиметрическая) карта Тихого океана и прилегающих территорий.}\label{fig:01}
\end{figure}

Двумя серьезными ограничивающими факторами исследования глубоководных желобов Мирового океана являются с одной стороны –  стоимость и трудоемкость организации морских экспедиций на основе НИС, и с другой стороны – сложность освоения скриптовой техники картографирования и обработки данных, которая требует сочетания междисциплинарных навыков (программирование, картография, геотектоника и геология). В связи с этим, данная диссертация является вкладом в существующие пробелы и нацелена обобщить информацию о глубоководных желобах Тихого океана путем системного картографирования на основе машинной обработки данных скриптами GMT.

\underline{\textbf{Вторая глава}} посвящена исследованию желобов Тихого океана.

Согласно классификации Международной гидрографической организации (IHO), глубоководные (океанические) желоба представляют собой «очень длинные, узкие, экстремально глубокие и асимметричные впадины морского дна с относительно крутыми сторонами». Глубоководные желоба представляют собой сложную систему со взаимосвязанными факторами: региональная геология, осадконакопление и мощность отложений, тектоника и геодинамика, сейсмическая активность. \\
Существуют различные геодинамические процессы, влияющие на динамику тектонических плит, структуру, а также геоморфологию рифтовых окраин. Образование глубоководных желобов вызвано погружением остывшей плиты в мантию и сопровождается землетрясениями и цунами. Геологическое строение Тихого океана, также как и его тектоническая структура имеют сложный характер. При этом, анализ батиметрии Тихого океана показывает более сложное и глубокое строение в его З и СЗ части. \\
Выбрано 20 глубоководных желобов в качестве объекта исследований в данном диссертационном проекте, которые включают следующие: 1) Алеутский, 2) Марианский, 3) Филиппинский, 4) Курило-Камчатский, 5) Центральноамериканский, 6) Перу-Чилийский, 7) Палау, 8) Японский, 9) Кермадек, 10) Тонга, 11) Идзу-Бонин, 12) Новобританский (Бугенвиль), 13) Сан-Кристобаль, 14) Манила, 15) Яп, 16) Новогебридский (Вануату), 17) Пюйсегюр, 18) Хикуранги, 19) Витязь и 20) Рюкю. Это основные желоба Тихого океана, наиболее представительные для геоморфологического моделирования. \\
	Морфоструктуры желобов разнообразны (V-, U-образная, дугообразная, удлиненная, перевернутая гребневидная, почти-овальная, симметричная, асимметричная: лево- или правосторонняя), что связано с преобладающими движениями тектонических плит и геологическим режимом и субстратом. Из-за специфических геологических условий океанические желоба имеют уникальное строение и особенности, сформированные в результате сложного процесса. Морфологически представлены как длинные, узкие, крутые впадины на дне океана, желоба расположены на максимальных глубинах, примерно от -7.300 до > -11.000 м. Существует ряд факторов со сложными взаимосвязями, влияющих на структуру и формирование желобов, в т.ч. геодинамика, геология, осадконакопление и тектоника. 

\underline{\textbf{Третья глава}} посвящена детальному описанию методологии и скриптовой техники визуализации и моделирования данных.

\underline{\textbf{GMT}} 

В процессе картографирования в рамках диссертации для визуализации карт всех желобов была применена картографическая программа Generic Mapping Tools (GMT), представляющая собой консольную программу, состоящую из отдельных модулей (свыше 100). В данной работе на основе существующих модулей GMT автором были разработаны новый методический подход для геоморфологического моделирования профилей в целях получения автоматических замеров продольно-поперечных профилей для получения массивов данных, содержащих данные глубин. Кроме того, автором разработана последовательность скриптов для 2D и 3D тематической визуализации геологических и геофизических карт 20 желобов Тихого океана. 

По сравнению с ГИС, преимущества GMT заключаются в его высокой эффективности благодаря использованию скриптового подхода в процессе картографирования, а не графического интерфейса пользователя (GUI), а также открытый исходный код. Общая схема методики GMT включала использование следующих модулей GMT для решения отдельных задач картографирования и пространственного моделирования желобов: 
\begin{enumerate}
	\item Модули 'grdcut' и 'grdimage' для визуализации GEBCO/ETOPO1; 
	\item Модуль “grd2cpt” для настроек цветовой палитры;
	\item Модули “grdconvert” и “img2grd” для конвертации форматов в GMT;
	\item Модуль “pscoast” для построения контуров;
	\item Модуль “psbasemap” для настроек карты (градусная сетка, подписи);
	\item Модуль “psxy” для картографирования точечных и символьных объектов;
	\item Использование модуля  “grdcontour” для построения изолиний;
	\item Использование модулей “pspstext”, для добавления аннотаций; 
	\item Трехмерное моделирование модулем “grdview”;
	\item Автоматизированная оцифровка ортогональных поперечных профилей модулем ’grdtrack’ вдоль выбранного сегмента желобов протяженностью 400 км, интервалом 20 км, с расстоянием 2 км между транзектами; 
	\item Модуль “'trend1d” для визуализации средних, медианных значений и трендов для профилей;
	\item Модули “pshistogram” и “psrose” для описательного статистического анализа (гистограмм);
	\item Модуль “nearneighbor” для аппроксимации поверхностей;
	\item Модуль “grd2xyz” для конвертации растра в табличный массив данных;
\end{enumerate}

В отличие от ГИС, GMT полностью основан на скриптовом подходе на основе UNIX с использованием командной строки без графического интерфейса. GMT совместим с GDAL, что было использовано для анализа метаданных и стат. информации о растрах (утилита “gdalinfo”). Функционал GMT намного превосходит традиционные ГИС. На основе многофункциональности GMT, скрипт использует многочисленные аргументы для тонкой настройки каждого картографического элемента: проекции, метки, добавление розы ветров, настройки главных и второстепенных сеток, заголовки и подзаголовки, разнообразие цветовых палитр и возможность подключать новые, тонкие настройки легенд, векторных элементов, аннотаций, и прочих элементов карты (береговые линии, масштабные линейки).

\underline{\textbf{AWK}} 
Следующий этапом работы было преобразование ряда таблиц, полученных в результате цифрования в GMT c помощью AWK. AWK является скриптовым языком, управляющим данными, запускаемым с консольной строки. AWK – мощный инструмент для извлечения данных, интегрированный с командами операционной системы Unix. AWK состоит из набора кодов-действий, что можно использовать в работе с потоками табличных данных. AWK позволяет извлекать или преобразовывать текст, производить форматированные записи, что и использовалось в настоящей диссертации при обработке большого массива данных для сокращения и преобразования таблиц. 

AWK использует строковый тип данных, обрабатывает ассоциативные массивы, индексированные ключевыми строками, и регулярные выражения. Обработанные геоданные были получены в численной форме в виде сложной многорядной таблицы для каждого сегмента ККЖ. Эти таблицы были созданы в GMT в его собственном формате и не могли быть обработаны напрямую в MATLAB/Octave. Поэтому таблицы были вначале экспортированы в AWK для автоматического форматирования (reshaping). Далее затем таблица была ре-структурирована, отсортирована и изменена скриптом AWK.

\underline{\textbf{Octave / MATLAB}} 

Экспериментальной новизной использования Octave в данной диссертации является интеграция этого статистического инструмента с пространственным изображением геоморфологических профилей желобов через AWK. Объединенные вместе для анализа геопространственных данных, эти инструменты позволили провести моделирование морфологии желобов на примере ККЖ. Поэтому основной целью данного раздела из цикла задач было представить разработанную техническую методологию для пакетной обработки пространственных данных с помощью цепочки пошагового рабочего процесса GMT-AWK-Octave. Данный технический алгоритм задач был использован для визуализации и анализа морфологической структуры желоба. 

GNU Octave является интерактивным языком высокого уровня с открытым исходным кодом, совместимым с MATLAB. После того, как таблицы были преобразованы в нужный формат с помощью манипуляций в AWK, они были обработаны скриптом Octave, в результате чего были визуализированы два графика, показывающих северный и южный сегменты ККЖ. Скрипт Octave был использован для визуализации и построения графиков и морфологического анализа профилей.

\underline{\textbf{Quantum ГИС}} 
Часть диссертационного проекта была выполнена в QGIS. Подводная геоморфология была изучена посредством пространственного и статистического анализа профилей поперечных сечений, оцифрованных в выбранных желобах (Марианский и Филиппинский). Обработка данных в QGIS включала различные аспекты обработки данных: поиск данных из веб-ресурсов, загрузка, структурирование и послойная визуализация данных на платформе QGIS, оцифровка и пространственная обработка с использованием плагинов, сбор и анализ данных и экспорт в таблицы csv, передача наборов данных в библиотеки R и Python, статистическая обработка, моделирование и визуализация графических графиков. 

Плагины QGIS разделяют несколько сложных задач на более мелкие и более управляемые компоненты: чтение координат, оцифровка поперечных профилей. В отличие от GMT, рабочий процесс оцифровки в QGIS был полуавтоматическим. Поперечные разрезы были оцифрованы в QGIS, их координаты и атрибуты значения сохранены в виде таблицы csv., далее использованные для статистического анализа в Python и R. 

\underline{\textbf{Python}} 
Использование Python, объектно-ориентированного языка программирования высокого уровня общего назначения,  является относительно новым подходом для задач тектоники и геодинамики,а также для анализа и обработки геоморфологических данных. В диссертации Python применяется для моделирования данных на основе применения встроенных математических алгоритмов и следующих статистических библиотек: Matplotlib и пр. Эффективность обработки и визуализации данных Python является ключевым фактором для применения его функциональности в стат. анализе двух выбранных желобов (Марианского и Филиппинского). 
	
Профили были проанализированы путем моделирования, направленного на анализ взаимосвязей между атрибутивными факторами с использованием основных существующих статистических методов. Статистический анализ проводился с использованием таблиц, содержащих атрибуты параметров оцифрованных желобов. Значения min/max и квартили были считаны в таблицах. 
	
\underline{\textbf{R}} 
Язык программирования R был использован для решения задач статистической обработки данных для двух выборочных желобов (Марианского и Филиппинского). R был выбран благодаря его логичному синтаксису и многофункциональным стандартным библиотекам. Используя пакеты R, была изучена неравномерность действия различных факторов, определяющих на морфоструктуру желобов (толщина отложений, крутизна склона, угловой аспект, глубины профиля и магматизм близлежащих территорий).

Статистическая анализ, выполняемый с помощью скриптов R, состоял из отдельных написанных кодов на основе алгоритмов машинной обработки табличных данных. Основные библиотеки R включали {ggplot2} для регрессионного анализа и построения кривых плотности распределения данных, диаграммы, {ggalt} для построения диаграмм для сравнения данных по тектоническим плитам, корреляционного анализа; {vcd} для мозаичных диаграмм, силуэтных участков для анализа композиционного сходства батиметрических профилей, ассоциативных участков; {car} для ANOVA и прочие. 


В \underline{\textbf{четвертой главе}} приведено описание результатов геоморфологического моделирования 20 глубоководных желобов Тихого Океана. Глубоководные океанические желоба представляют собой важную часть океанической литосферы: крупнейшие по протяженности, линейные углубленные элементы рельефа поверхности морского дна Земли, отличающиеся своей глубиной, непрерывностью и уникальной морфологией склонов. Желоба преимущественно расположены вдоль окраин Тихого океана. В основном, они развиты вблизи островных дуг и орогенов, образующихся над зонами субдукции. Например, максимально длинный Перу-Чилийской желоб имеет длину 4.500 км и достигает глубин 7–8 км под уровнем моря. 

Желоба, расположенные на западе Тихого океана, глубже, чем на восточных его границах. При этом, наибольшие глубины желобов достигают до 10–11 км, обнаруженные в желобах Марианском и Тонга-Кермадек. В настоящей работе были определены следующие типы формы морфологических профилей, и желоба классифицированы на следующие типы: U-образные, V-образные, асимметричные левосторонние, серповидные, извилистые гребневидные, каскадные, округло-дугообразные, асимметричные правосторонние. Для каждого геоморфологического типа характерны подтипы крутизны, определены: сильные, очень сильные, экстремальные, крутые, очень крутые склоны. 

Кроме того, склоны долин желобов классифицируются как очень высокие, высокие, средние, низкие по степени кривизны. Следовательно, классы размеров и классы склонов долин анализируются в контексте физической среды, а также тектонических и геологических условий участка желоба. Форма изгиба погружающейся литосферной плиты также оказывает влияние на тектоническую нагрузку на оси желоба, прочность, а также на разломные эффекты погружения тектонической плиты, что отражается в разнообразии форм желоба. 

\pdfbookmark{Курило-Камчатский желоб}{kkt}  
\underline{\textbf{Курило-Камчатский желоб}}. Курило-Камчатская вулканическая дуга и сопряженный желоб связаны с динамикой Камчатско-Алеутского сочленения и субдукцией Тихоокеанской плиты. Регион отличает повышенная сейсмичность, вулканизм и землетрясения, свидетельствующие о тектонической активности региона, что обуславливает особые геохимические условия донных отложений. Вариативность геоморфологии KKT заключаются в особых тектонических условиях на западе Тихого океана. 

\begin{figure}[ht]
    \centerfloat{
        \hfill
        \subcaptionbox{Батиметрия}{%
            \includegraphics[width=0.47\linewidth]{Рис_02a.jpg}}
        \hfill
        \subcaptionbox{Профили}{%
            \includegraphics[width=0.53\linewidth]{Рис_02b.jpg}}
        \hfill
    }
    \caption{Курило-Камчатский желоб}\label{fig:02}
\end{figure}

Курило-Камчатский желоб (ККЖ): желоб каскадной формы (Рис. \ref{fig:02}). Южная часть достигает максимальных глубин -8.200 м, а северная – 7.800 м. Морфология желоба на севере окончания южного сегмента ККЖ при подходе к проливу Буссоль имеет U-образную форму. Клиновидная форма желоба выражена в его южной части на западном склоне. На участке вдоль северной части желоба определены особенности геоморфологической системы (подводные каньоны, клин желоба, подводные гребни и котловины). Клинообразная или U-образная морфология желоба в двух частях (южный и северный) отличается от внутреннего склона и в продольном направлении по оси желоба. Профили южного сегмента желоба, по сравнению с северным сегментом, расположены вблизи интенсивной вулканической зоны, вследствие чего в этой части присутствуют эрозионные рвы на глубине -3000 м. 

Ортогональные батиметрические профили через южный и северный сегменты KKЖ идентифицируют две различных U- и V-образных морфологических особенностей с воздействием на доминирующие процессы: (1) в северном сегменте доминируют процессы отложения наносов, поступающие с п-ова Камчатка; (2) южный сегмент отличается процессами эрозии и интенсивным переносом наносов. Северная часть показывает более глубокие значения и менее крутые склоны по сравнению с к южному сегменту, который относится к размывающим рвам и дрейфу подводных отложений. Южная часть имеет большие глубины и более крутой уклон, что можно объяснить интенсивной сейсмичностью и вулканизмом. ККЖ образует V-образную долину в своем южном сегменте с шириной профиля около 3 км, и корытообразную форму (П-образную в перевернутой форме) в северном сегменте.

\pdfbookmark{Алеутский желоб}{at}  
\underline{\textbf{Алеутский желоб}}. Алеутский желоб (АЖ) имеет удлиненную корытообразную форму (Рис. \ref{fig:03}). Коллизия Кроноцкой дуги на северо-восточной окраине Евразии значительно повлияла на структурную эволюцию Камчатско-Алеутского сочленения, а тажке распределение очагов землетрясений. Земная кора имеет переходную структуру от океанической к континентальной на восточной части Алеутских островов и в районе п-ва Аляска. Северо-западная часть АЖ со стороны цепи островов имеет более крутой и короткий уклон с расположенным рядом морских террас. Средний уклон его северной стороны колеблется от 3° до 4°, в то время как на его более мелкой южной стороне - 1° или 2°. Согласно моделированию статистической гистограммы, наиболее частое значение глубин в исследуемой области желоба – 4.800 с частотой 1.722. 

\begin{figure}[ht]
    \centerfloat{
        \hfill
        \subcaptionbox{Алеутский желоб}{%
            \includegraphics[width=0.52\linewidth]{Рис_03a.jpg}}
        \hfill
        \subcaptionbox{Центральноамериканский желоб}{%
            \includegraphics[width=0.47\linewidth]{Рис_03b.jpg}}
        \hfill
    }
    \caption{Алеутский и Центральноамериканский желоба}\label{fig:03}
\end{figure}

Статистический анализ данных показал, что глубины с диапазонами от -5600 – -5400 м отмечены в 1106 данных (8\%); глубины с диапазонами -5400 до -5200 м имеют повторяемость 846 проб (7\%); глубины с диапазонами -4900 и -5200 м имеют повторяемость 1.297 (11\%). Самые глубокие значения зафиксированы как -4.900 и -4.700 при 1.722 пробах (13,2\%). Более мелкие значения продемонстрировали мелководные участки с диапазоном глубин от -4500 до -3000 м с промерами менее 200 ( <2,5\%). Значения <-500 м, относятся к шельфовым районам вблизи Алеутских о-вов.

\pdfbookmark{Центральноамериканский желоб}{mat}  
\underline{\textbf{Центральноамериканский желоб}}. Центральноамериканский желоб (ЦЖ) имеет симметричную форму на Гватемальском отрезке (Рис. \ref{fig:03}). Желоб простирается вдоль западного берега Центральной Америки от центральной части Мексики до Коста-Рики. Морфологическое строение имеет прямую форму склонов с выражено более крутой передней дугой к стороне океана. Желоб образовался в результате субдукции Тихоокеанской плиты. Центральноамериканский желоб можно условно разделить на северную (желоб Акапулько) и южную (Гватемальский желоб) части. Желоб Акапулько проходит от Халиско до хр. Теуантепек, а Гватемальский желоб, проходит от хр. Теуантепек до хр. Кокос.

Геометрия ЦЖ имеет крутую и прямую форму. Основные выборки глубинных наблюдений имеют значения от -3000 до -6200 м. Дно желоба имеет среднюю ширину 3-5 км. Есть различия в склоне желоба, обращенных к океану и более крутой склон на материковой стороне. В других областях террасы менее четко очерчены. По сравнению с другими желобами Тихого океана, морфология ЦЖ представляет выпрямленные крутые склоны без значительной кривизны форм. Пространственный анализ, выполненный с помощью GMT, показал, что основные точки глубинных наблюдений ЦЖ расположены между -3000 и -6200 м. Самый глубокий диапазон данных батиметрии с 1.111 записями составляет 25\% от общего пула точек промеров. Средние значения выборок наблюдений имеют значение -2200, в то время как медианные значения имеют значения -3200 до -3300 м.4.4  

\begin{figure}[ht]
    \centerfloat{
        \hfill
        \subcaptionbox{Сегменты желоба}{%
            \includegraphics[width=0.3\linewidth]{Рис_04a.jpg}}
        \hfill
        \subcaptionbox{Модели профилей}{%
            \includegraphics[width=0.55\linewidth]{Рис_04b.jpg}}
        \hfill
    }
    \caption{Перу-Чилийский желоб}\label{fig:04}
\end{figure}

\pdfbookmark{Перу-Чилийский желоб}{pct}  
\underline{\textbf{Перу-Чилийский желоб}}. Перу-Чилийский желоб (ПЧЖ) имеет асимметричную форму и структурно состоит из двух частей (Рис. \ref{fig:04}). Желоб является результатом конвергентной границы тектонических плит, где восточный край плиты Наска погружается под континентальную Южно-Американскую плиту. Расположенный в восточной части Тихого океана у берегов Перу и Чили, это самая длинный по географической протяженности глубоководный желоб в мире, неоднократно рассматриваемый в литературе.
Чилийский сегмент отличается большими глубинами и крутым уклоном на плите Наска, в то время как Перуанский сегмент желоба демонстрирует более пологий рельеф со стороны океана. В качестве сравнения наиболее повторяющихся  данных для обоих сегментов можно привести следующие: перуанский сегмент имеет большинство данных (23\%), 1410 образцов на глубинах от -4500 м до -4700 м. Он показывает пик всех данных с очень крутая картиной в распределении. Остальные данные в соседнем диапазоне намного ниже: 559 выборок на глубинах от -4700 до -5000 м и 807 выборок на глубинах от -4200 до -4400 м для перуанского сегмента. В чилийском сегменте более равномерное распределение данных. Кроме того, диапазон данных составляет 205 проб для чилийского сегмента на глубинах -4700 до -5000 м (против 559 для перуанского сегмента); и 731 для чилийского сегмента на глубинах от -4200 м до -4400 м (против 807 на перуанском участке); 317 для чилийского сегмента на глубинах от -4500 м до -4700 м (против 1410 в перуанском сегменте). Наиболее частая глубина для перуанского сегмента ПЧЖ - это диапазон от -4000 м до -4200 м (827 проб) против диапазона -4500 м до -4700 м для перуанского сегмента (1410 выборок). Перуанский сегмент в целом глубже и круче и характеризуется крутыми спусками. 

\pdfbookmark{Кермадек желоб}{kt}  
\underline{\textbf{Желоб Кермадек}}. Желоб Кермадек (ЖК): вытянутый желоб по перевернутой гребневидной форме (Рис. \ref{fig:05}). Желоб Кермадек расположен в южной части Тихого океана. Он простирается на ок. 1000 км от цепи подводных гор Луисвилл на севере до плато Хикуранги на юге к СВ от о. Северный в Новой Зеландии.  Вместе с желобом Тонга на севере он образует почти 2000 км линейную систему субдукции Кермадек-Тонга с особыми геохимическими условиями в результате вулканизма и магматизма. Субдукция начала развиваться в палеогене (эоцен) при погружении Тихоокеанской плиты под Австралийскую. Профили оцифрованы перпендикулярны оси желоба. 

\begin{figure}[ht]
    \centerfloat{
            \includegraphics[width=0.75\linewidth]{Рис_05.jpg}}
    \caption{Желоб Кермадек и Тонга.}\label{fig:05}
\end{figure}

Широта возле каждого профиля указывает на положение осевого дна. ЖК имеет плавную форму западного склона и резкую форму на восточном фланге с глубинами от -6600 до -6800 м согласно анализу 2641 промеров глубины. Сравнивая самые глубокие значения (свыше -9000 м) для ЖК, согласно анализу 251 образцов, очевидно, что ЖТ глубже, чем ЖК. Для глубин в диапазоне от -6000 до -5000 м больше значений: 6167 выборок. ЖК имеет пологий склон и небольшую глубину. Это иллюстрирует тектонические и геологические локальные вариации, а также вариативность процессов осадконакопления в ЖК, вызывающая вариации их морфоструктуры, несмотря на их близкое расположение. 

\pdfbookmark{Тонга желоб}{tt} 
\underline{\textbf{Желоб Тонга}} 
Желоб Тонга продолжает субдукционную систему Кермадек на севере, примерно от 24° ю.ш., где Тихоокеанская плита быстро погружается под Индо-Австралийскую вдоль желоба Тонга-Кермадек в ЮЗ части океана. 
\begin{wrapfigure}{r}{0.5\textwidth}
  \begin{center}
            \includegraphics[width=0.9\linewidth]{Рис_06.jpg}
                \end{center}
   \caption{Желоб Кермадек и Тонга: профили}\label{fig:06}
\end{wrapfigure}
Базальтовые лавы, извлеченные из верхней коры входящей литосферы по длине желоба Тонга, имеют средне- и позднемеловой возраст, в основном представленyst N-MORB, варьируются от толеитовых до щелочных базальтов. Самая активная зона мантийной сейсмичности Земли возникает в результате субдукции Тихоокеанской плиты в желобе Тонга, причем до сих пор неизвестно, почему эта плита вызывает намного больше землетрясений, чем другие погружающиеся плиты в мире. Желоб Тонга (ЖТ) имеет меньшую глубину по сравнению с ЖК (Рис. \ref{fig:05}), несмотря на близкое расположение. Сравнение самых глубоких значений (свыше -9000 м) показывает для ЖТ 320 отборов проб на данных глубинах. Сравнивая варианты глубин морского дна в диапазоне от -6000 до -5000 м, можно отметить, что ЖТ имеет 4803 выборок. Автоматическая оцифровка профилей поперечного сечения показали, что геоморфология ЖТ имеет более крутой уклон на западном фланге. Результаты показывают большие значения глубин, неровность дна и крутой уклон профилей для ЖT по сравнению с ЖК (Рис. \ref{fig:06}).

\pdfbookmark{Вануату желоб}{vt} 
\underline{\textbf{Желоб Вануату (Новые Гебриды)}}, (Рис. \ref{fig:07}). 
Желоб Ваунату расположен в ЮЗ части Тихого океана, на окраине Кораллового моря, между о. Вануату и Новая Каледония. Он превышет глубины 7600 м. Желоб составляет около 1200 км в длину и 70 км в ширину, огибает южную оконечность Вануату, образуя диагональную симметрию с северным концом близлежащего желоба Тонга. Вулканическая островная дуга являестя источником магмы и вулканизма в регионе.

\begin{figure}[ht]
    \centerfloat{
        \hfill
        \subcaptionbox{Сегменты желобов}{%
            \includegraphics[width=0.47\linewidth]{Рис_07a.jpg}}
        \hfill
        \subcaptionbox{Модели профилей}{%
            \includegraphics[width=0.53\linewidth]{Рис_07b.jpg}}
        \hfill
    }
    \caption{Желоба Вануату и Витязь}\label{fig:07}
\end{figure} 

Результаты моделирования профилей поперечного сечения для желоба Вануату (ЖВ) продемонстрировали, что ЖВ достигает значений -7000 м). Форма гистограммы различается двумя полосами: ЖВит имеет распределение данных в виде колокола с двумя пиками при значениях глубин от -3500 до -3700 м, а второй — на диапазоне от -4800 до -5000 м, то есть большинство значений находятся в общем диапазоне от -5000 до -3000 м. Здесь отмечено 1605 отборов (517, 653 и 435). Из них первый пик имеет значения глубин от -3300 до -4000 и 839  наблюдений (238, 304, 297) и далее на глубинах от -4500 до -5200 м. Окружающий рельеф двух желобов варьируется: для ЖВит прилегающая абиссальная равнина имеет выположенный рельеф, в то время как для ЖВан окружен более сложным подводным рельефом бассейна северного Фиджи (восточный фланг), так и западного побережья.

\pdfbookmark{Витязь желоб}{vtt} 
\underline{\textbf{Желоб Витязь}}. Результаты моделирования ЖВит показывает глубины не превышающие -6.100 м (Рис. \ref{fig:07}). По сравнению с Жвит, ЖВан имеет более V-образную классическую форму для желоба с пологими формами на обоих склонах (западный и восточный). Сравнение средних значений двух желобов показывает, что выбранный участок ЖВит имеет меньшие глубины с максимальным 5.000 м, в то время как ЖВан имеют глубины с -6.000 м. Самые глубокие значения (свыше -5000 м) для ЖВит отмечены для 413 промеров. Морфоструктура  двух желобов существенно различается: ЖВит имеет более плоское широкое корытообразное дно, с более крутым уклоном на восточном фланге. Сравнительный анализ геофизических сеток выявил структурные характеристики двух океанических желобов, ЖВан и ЖВит, и прилегающей океанической коры северной части бассейна Фиджи. 

Согласованность растров с данными аномалии силы тяжести в редукции Фая и вертикальным градиентом силы тяжести подтверждают, что данные совпадают для 10 мГал, диапазон 3-10 мГал означает переходную зону гравитационной аномалии, в которой увеличивается когерентность. В интервале 10-20 мГал когерентность снижается до очень узкого интервала, и корреляция совпадает до 100 мГал. Согласованность гравитационного градиента добавляет новую информацию к интерпретации данных, охватывающих активные сходящиеся границы в Тихом океане. Распределение данных для ЖВан при этом более равномерно при небольшом увеличении глубиной (от -2200 до -800 м), для 1353 проб, что отмечено на соответствующих интервалах на гистограмме, а также на карте батиметрического моделирования данных

\pdfbookmark{Новобританский желоб}{nbt} 
\underline{\textbf{Желоб Новобританский (Бугенвиль)}}. Новобританский желоб (НБЖ) имеет асимметричное U-образное поперечное сечение профиля, с выраженным уклоном на запад в восточном направлении между 35° и 41° (Рис. \ref{fig:08}). Желоб географически расположен в Соломоновом море, расположенном в ЮЗ Тихого океана, между Папуа-Новой Гвинеей и Соломоновыми островами. Бассейн Вудларк, расположенный к югу от дуги Соломоновых островов, представляет собой молодой (~ 5 млн лет назад) океанический бассейн, который погружается под Ново-Британcкий желоб Этот регион является одной из немногих зон субдукции в мире, где происходит субдукция молодой плиты длительностью несколько миллионов лет. Увеличение глубины НБЖ имеет постепенный плавный характер: на отрезке 50 км достигая -2000 м на фланге со стороны Соломонова моря. 

\begin{figure}[ht]
    \centerfloat{
        \hfill
        \subcaptionbox{Сегменты желобов}{%
            \includegraphics[width=0.53\linewidth]{Рис_08a.jpg}}
        \hfill
        \subcaptionbox{Модели профилей}{%
            \includegraphics[width=0.47\linewidth]{Рис_08b.jpg}}
        \hfill
    }
    \caption{Желоба Новобританский (Бугенвиль) и Сан-Кристобаль}\label{fig:08}
\end{figure}

На графике поперечного сечения НБЖ заметно асимметричен и имеет наклон в сторону суши, дугообразный на востоке. НБЖ более мелкий по сравнению с ЖСК имея глубины в среднем <-7000 м. Самые высокие значения выборки регистрируются в диапазоне глубин от -5000 до -5200 м, (272 образцов), за которыми постепенно следует диапазон глубин от -5200 до -5000, далее от -5000 до -4800 и от -4.800 до -4.600 м (кол-во образцов: 204, 189, 157). Глубины НБЖ постепенно увеличиваются в сегменте от -50 до 0 м на профиле в поперечного сечения желоба, и уменьшаются на диапазоне от -3.000 до -7.000 в западном направлении, достигая максимальных глубин при -7.000 м и среднего значения при -8.000 м. Гистограмма имеет усеченный вид по правому флангу желоба с преобладающими значениями на малых глубинах и уменьшающимися <-5.000 m по сравнению нормальным распределением данных ЖСК. 

\pdfbookmark{Сан-Кристобаль желоб}{sct} 
\underline{\textbf{Желоб Сан-Кристобаль}}. Желоб Сан-Кристобаль (ЖСК) симметричен и имеет V-образную форму для отрезка от -50 до -50 км на графике (Рис. \ref{fig:08}). ЖСК имеет градиент склона 33° на склоне, обращенном в сторону океана и 33,69° континентального склона, при этом характер рельефа более прямой на профиле со стороны океана. Подъем ЖСК более резкий по профилю: на расстоянии 0-50 км глубины поднимаются от -8000 до -5200 м. Для участка разреза на расстоянии -50 до 0 м (центральной оси желоба) глубины становятся больше, после чего уменьшаются от -3000 до -8000 м в западном направлении, то есть наклон становится круче. Морфоструктура объясняется контекстом исторических и современных геологических и осадочных процессов, которые в совокупности повлияли на его формирование. На океанской стороне выделены несколько отдельных мелких каньонов и борозд.

\pdfbookmark{Филиппинский желоб}{pt} 
\underline{\textbf{Филиппинский желоб}}. 

Филиппинский желоб (ФЖ), один из глубочайших в мире, расположен на оси столкновения плиты Филиппинского моря (ПФМ) и плиты Сунда, западной части Тихого океана в Филиппинском море (ФМ) (Рис. \ref{fig:09}). Рельеф дна ФМ имеет черты переходной зоны с крупными геоморфологическими формами и осадконакоплением. Строение котловины моря сложное: здесь отмечены подводные горы, хребты большой протяженности, подводные вулканы, гайоты. Морфология ФЖ отмечена сложным взаимодействием трех литосферных плит: Евразийской, Австралийской и ПФМ, который включает столкновение, субдукцию и аккрецию. Древняя, тяжелая и плотная (103 300 000 $км^2$) Тихоокеанская плита погружается под Австралийскую (47,000,000 $км^2$) и ПФМ 5,500,000 $км^2$. ФЖ имеет классическую колоколообразную форму распределения глубин, показывающих заметное резкое увеличение частоты данных на интервале от -5.000 до -6.000 м ниже уровня моря: 1.117 выборок ( 37\% данных). Желоб имеет V-образную форму в самых глубоких сегментах. 

\begin{wrapfigure}{r}{0.5\textwidth}
  \begin{center}
            \includegraphics[width=0.9\linewidth]{Рис_09.jpg}
    \end{center}
   \caption{Филиппинский желоб}\label{fig:09}
\end{wrapfigure}

График показывает изменение углов склона желоба, крутизны и перепад высот с западного и восточного флангов желобов. 3D-модель показывает выбранный сегмент ФЖ, где край характеризуется крутым уклоном и примыкает к о. Минданао, к северу от отложений наносов и к СЗ из-за тектонической эрозии. 

\pdfbookmark{Манильский желоб}{mt} 
\underline{\textbf{Манильский желоб}}. 
Манильский желоб, расположенный к западу от о. Лусон и Миндоро на Филиппинах, достигает глубины ок. 5400 м, в отличие от средней глубины Южно-Китайского моря около 1500 м. Он создается субдукцией, при которой плита Сунда (часть Евразийской плиты) погружается под Филиппинскую, образуя желоб с почти меридиональным простиранием в регионе активных тектонических движений, в т.ч. частых цунами и землетрясений. Конвергентная граница заканчивается на севере зоной столкновения Тайвань, а на юге террейном Миндоро (блок Сулу-Палаван, сталкивающийся с юго-западом Лусона).

\begin{figure}[ht]
    \centerfloat{
        \hfill
        \subcaptionbox{Сегменты желоба}{%
            \includegraphics[width=0.55\linewidth]{Рис_10a.jpg}}
        \hfill
        \subcaptionbox{Гравиметрия в редукции Фая}{%
            \includegraphics[width=0.45\linewidth]{Рис_10b.jpg}}
        \hfill
    }
    \caption{Манильский желоб}\label{fig:10}
\end{figure}

Результаты геоморфологического моделирования и картографирования Манильского желоба (МанЖ) показывают пространственную вариативность в разных частях желоба (Рис. \ref{fig:10}). Проведена характеристика геологии морского дна путем анализа положения желоба и тектонических плит, а также площадей распространения подводных вулканов. Сопоставление профилей разрезов демонстрирует различия морфоструктуры в южной и северной частях желоба. Выявлено, что южная часть имеет более крутой уклон в сторону океана, тогда как северная часть круче на континентальной части склона. Подводные террасы расположены в северной части желоба на глубине -2000 м.

Геофизическое моделирование, основанного на данных аномалий силы тяжести в редукции Фая, показывает точки высокой плотности распространяются непрерывно с севера на юг Тайваня вертикально на юг до о. Лусон, Филиппины, от -50 мГал до -200 мГал. Большинство значений в районе Южно-Китайского моря находится в диапазоне от 10 до 30 мГал, тогда как значения свыше 60 мГал могут быть отмечены в островных районах и южных Филиппинах. Результаты геоморфологического моделирования и картографирования МанЖ выявили основные тектонические условия, в которых столкновение плит вызывает нестабильность геологических условий, влияющих на формирование морфоструктуры МанЖ: повторяющиеся землетрясения, подводный вулканизм и повышенную сейсмичность.

\pdfbookmark{Палау желоб}{plt} 
\underline{\textbf{Желоб Палау}}. 
ЖПал входит в систему желобов Палау и Яп в СЗ части Тихого океана и обладает уникальными морфологическими и геологическими характеристиками: достигает глубин свыше 8000 м, несмотря на чрезвычайно малую скорость конвергенции плит и относительно молодой возраст субдуцированной плиты; V-образное поперечное сечение желобов, объясняет небольшой уровень осадконакоплений; крутые склоны желобов в направлении к суше; аномально близкое расположение о-вов Палау и Яп к оси желоба; преобладающее распределение барьерных рифов на западной периферии, противоположное к доминантному направлению ветров с В на СВ.

Результаты анализа данных продемонстрировали пространственную морфологическую неравномерность желоба Палау (Рис. \ref{fig:11}). Ранжирование геологических данных включает анализ степеней градиентов склонов, которые были определены пространственным моделированием.Вариативность морфоструктур  основан на распределении данных профилей ЖПал, который имеет характерную закономерность: количество глубин в диапазоне от -5000 м до -4000 м отчетливо выделяется по сравнению с диапазонами на других глубинах (737 проб). Такая же ситуация наблюдается для диапазона глубин от -5000 до -6000 м, где насчитано 152 образца. В целом, ЖПал имеет пологий уклон как к суше, так и к океану. Профили оси ЖПал имеют V-образную форму и пологий склон в направлении к суше и океану. Контуры аномалий силы тяжести в открытом воздухе в редукции Фая в целом коррелируют с изолиниями. 

\begin{figure}[ht]
    \centerfloat{
        \hfill
        \subcaptionbox{Сегменты желобов}{%
            \includegraphics[width=0.45\linewidth]{Рис_11a.jpg}}
        \hfill
        \subcaptionbox{Модели профилей}{%
            \includegraphics[width=0.55\linewidth]{Рис_11b.jpg}}
        \hfill
    }
    \caption{Желоба Яп и Палау}\label{fig:11}
\end{figure}

\pdfbookmark{Яп желоб}{yt} 
\underline{\textbf{Желоб Яп}}. 
Желоб Яп продолжает региональную систему желобов Яп-Палау в западной части Тихого океана, представляя собой эволюцию мантии под зоной субдукции, что выражено в особом геохимическом составе: серпентинизированные перидотиты, габбро, базальты и известняки обнажены на склоне, обращенном к суше, без осадочного чехла и покрытий из гидрооксида ферромарганца. Сравнение перидотитов из деплетированной мантии, предполагает, что мантия желоба Яп была модифицирована субдукционным расплавом, характеризующимся высоким содержанием и соотношениями Zr/Hf. Плавление могло быть усилено тектонической эрозией погружающейся плиты Кэролайн со сложными тектоническими морфоструктурами на самых ранних стадиях инициирования субдукции.

Согласно статистическому анализу, наибольшие значения желоба Яп (ЖЯ) составляют от -4000 м до -4200 м (278 выборок) с 19\% и 17\% от общего набора данных, соответственно (Рис. \ref{fig:11}). Распределение данных для ЖЯ (605 выборок) следующее: диапазон глубин от -5000 м до -6000 м: в ЖЯ зарегистрировано 53 образца. В целом, ЖЯ имеет более крутые склоны и большие глубины наблюдений, а также более глубокие значения в диапазоне от -7000 м до -6000 м: 68 проб, в то время как соседний желоб Палау имеет только 34 точки для того же диапазона глубин [261]. В целом, ЖЯ имеет более крутые склоны и большие величины глубин наблюдений. ЖЯ также характеризуется наличием тонкого осадочного чехла на дне желоба по его оси, что доказывает взаимосвязь между внешними факторами (геология, геоморфология, геофизика, тектоника, седиментация), влияющими на морфоструктуру и динамику желоба. Топографические профили оси ЖЯ имеют очень четкую V-образную форму, что указывает на практически полное отсутствие осадков на дне желоба либо очень тонкий осадочным покров.

\pdfbookmark{Рюкю желоб}{rt} 
\underline{\textbf{Желоб Рюкю}}. 
Желоб Рюкю (ЖР) представляет собой океанический желоб длиной 1400 км, протянувшийся вдоль юго-восточного края островной дуги Рюкю. Географически, желоб Рюкю расположен в субтропической регионе Тихого океана между Филиппинским морем, о. Тайвань и о. Кюсю, южная Япония (Рис. \ref{fig:12}). Желоб образовался в результате погружения океанической коры плиты Филиппинского моря под континентальную кору Евразийской плиты. Исследования показали различия в морфологии северных и южных сегментов, что отражает корреляцию между формой желоба, крутизной его градиентных углов откоса и локальными тектоникой региона. В северной части желоба субдукция плиты Филиппинского моря отличается небольшой глубиной и крутизной, достигающей ок. 11°. Напротив, угол становится круче, достигая 70° ниже ок. 70 км. Регион, расположенный к  СЗ от вулканического острова Рюкю - это задуговый бассейн впадины Окинавского прогиба. 

\begin{figure}[ht]
    \centerfloat{
        \hfill
        \subcaptionbox{Геологическая обстановка}{%
            \includegraphics[width=0.40\linewidth]{Рис_12a.jpg}}
        \hfill
        \subcaptionbox{Модели профилей}{%
            \includegraphics[width=0.55\linewidth]{Рис_12b.jpg}}
        \hfill
    }
    \caption{Желоб Рюкю}\label{fig:12}
\end{figure}

Статистические результаты моделирования ЖР показывают различия в диапазонах глубин на разных его сегментах вдоль островной дуги Рюкю. Глубина и геометрия склона сильно различаются: максимальная значения от -5500 до -5700 м примечательны для северной части желоба, где отмечены 334 точек наблюдений. С небольшой разницей, южная часть демонстрирует сопоставимые значения (329 образцов) для того же диапазона. Более значительные различия видны для глубин свыше -6000 м. Южная часть желоба имеет лишь несколько значений (9 образцов) с глубинами свыше -6800 м, тогда как северная показывает 41, 59 и 23 измерения с максимальными значениями, достигающими -7460 м. 

Диапазон -3000 до -5000 м, включающий от умеренных до небольших глубин в сопоставимых для обеих частей желоба, значения не превышают 68 проб для северной части и 42 пробы для южной, соответственно. В преддуговом бассейне глубиной более 2000 м и менее 1000 м под уровнем моря показано больше данных проб в северной части желоба, чем в южной. Напротив, широкая терраса с глубиной менее 25 м показывает больше проб (161) в южной части желоба. Сравнение между двумя пространственно разными сегментами желоба показывают, что северная часть желоба имеет более глубокие отметки глубин по сравнению с южными. Это также можно отметить в глубоководной зоне, простирающейся в сторону моря от разлома Мияко, где зафиксированы самые глубокие значения. 

\pdfbookmark{Марианский желоб}{mt} 
\underline{\textbf{Марианский желоб}}. Географическое положение самой глубокой точки Мирового океана с максимальной глубиной 10.984 м. расположено в MрЖ, на глубине 11,329903°с.ш./142,199305°в.д. на глубине Челленджера. Как глубочайшая точка Земли, MрЖ неуклонно привлекает внимание научного мира (Рис. \ref{fig:13}). Серпообразная дуга MрЖ может быть уложена в прямоугольник с координатами 140°E-148°E, 11°N-25°N. Батиметрия морского дна варьируется в сторону моря от оси желоба: область Челленджера, к юго-западу от о. Гуам, ассоциируется с глубоким рельефом 10.984,25 м (95\%) при 11,329903°N/142,199305°E. 

\begin{figure}[ht]
    \centerfloat{
        \hfill
        \subcaptionbox{Сейсмическая обстановка}{%
            \includegraphics[width=0.45\linewidth]{Рис_13a.jpg}}
        \hfill
        \subcaptionbox{Модель геоида}{%
            \includegraphics[width=0.53\linewidth]{Рис_13b.jpg}}
        \hfill
    }
    \caption{Марианский желоб}\label{fig:13}
\end{figure}
МрЖ является частью субдукционной системы Идзу–Бонин–Мариана, которая геометрически представляет собой дугообразную форму полумесяца, и формирует границу и взаимоотношение между двумя тектоническими плитами: Тихоокеанской, субдуцирующей под Филиппинскую. Cложное морфологическое строение и неравномерность профилей MрЖ отражает его геолого-тектоническую эволюцию. Асимметрия строения проявляется на ЮЗ, где морфологическая структура желоба меняется и расположена глубочайшая точка на Земле. Морфология подвержена воздействию различных факторов, которые влияют на его форму (Рис. \ref{fig:14}). К ним относятся батиметрия (глубины), субдукция тектонических плит, геология, процессы осадконакопления и геометрия (крутизна склона), что рассмотрено в литературе. Границы субдукции Тихоокеанской плиты под Филиппинскую классифицируются как «Марианский тип», где связь между двумя плитами либо слабая, либо отсутствует. Западный склон MрЖ представляет собой откос без поднятия, а восточная часть является границей, вдоль которой проходит западный край Тихоокеанской плиты, погруженный под меньшую по размеру Марианскую плиту. 

\begin{wrapfigure}{r}{0.5\textwidth}
  \begin{center}
            \includegraphics[width=0.9\linewidth]{Рис_14.jpg}
    \end{center}
   \caption{Марианский желоб}\label{fig:14}
\end{wrapfigure}

В результате геоморфологических особенностей MрЖ сформированы его очень крутые склоны и глубины 3-5 км, ширина днища редко превышает 5 км. MрЖ имеет сложные ступени различной формы и размера, вызванные активными тектоническими процессами и эрозией. Крутизна его склонов в среднем составляет 4-5°, но отдельные участки могут ограничиваться более крутыми склонами, поскольку они подвержены гравитационному потоку в подводных каньонах и долинах. Сложное распределение разнообразного материала на прилегающих абиссальных равнинах способствует формированию морфологических особенностей регион MрЖ. Основная часть дна желоба имеет преобладающие глубины свыше 3000 м. Дно Тихого океана в районе МрЖ представляет собой самую древнюю его часть, образовавшуюся в поздней юре. 

Районы с подводными горами на оси желоба отмечены его пологой геоморфологией, меньшими осевыми вертикальными нагрузками и выпуклым рельефом на внешнем подъеме. Марианская дуга и Западно-Марианский хребет граничат с желобом с западной стороны. Каролинский хребет и группа Каролинских островов расположены к югу от оси желоба на расстоянии около 0-250 км. МарЖ имеет более крутой уклон на северном фланге со стороны архипелага. Крутизна склона на левом фланге составляет 33,10° в 0-15 км. Сегмент 15-25 км имеет крутизну 45°, сегмент 25-40 км - 30,75°, сегмент 40-50 км - 28°. На правом фланге желоба крутизна его градиента меняется следующим образом. Сегмент на дистанции 0-28 км имеет крутизну 40,5° от отвеса, 28-37 км от точки старта имеет 43,33° крутизны. Затем он постепенно становится пологим до 30,43° до 50 км. Сегмент на 50-90 км заключается в дугообразном постепенном повышении абсолютных значений глубин, которое стабилизируется через 100 км на уровне менее -4.000 м.

\pdfbookmark{Идзу-Бонин желоб}{ibt} 
\underline{\textbf{Желоб Идзу-Бонин}}. Желоб Идзу-Бонин (ИБ), также известный как Идзу-Огасавара, расположен на западе Тихий океан к югу от Японии (Рис. \ref{fig:15}). В продольном направлении он простирается от г. Токио, Япония, до Гуама. По своему строению ИБ продолжает Японский желоб. ИБ сформирован в зоне активной тектонической субдукции мелового периода Тихоокеанской плиты под плиту Филиппинского моря. Как следствие субдукции литосферных плит формируются следующие структуры: вулканическая дуга островов Идзу и Бонин, дуговая система Идзу-Бонин-Мариана и бассейн Сикоку как тыловой бассейн. С востока к ИБ примыкает плато Огасавара.

\begin{figure}[ht]
    \centerfloat{
        \hfill
        \subcaptionbox{Батиметрия}{%
            \includegraphics[width=0.52\linewidth]{Рис_15a.jpg}}
        \hfill
        \subcaptionbox{Тектонико-геологические условия}{%
            \includegraphics[width=0.48\linewidth]{Рис_15b.jpg}}
        \hfill
    }
    \caption{Желоб Идзу-Бонин}\label{fig:15}
\end{figure}

\begin{figure}[ht]
    \centerfloat{
        \hfill
        \subcaptionbox{Сегменты желоба}{%
            \includegraphics[width=0.45\linewidth]{Рис_16a.jpg}}
        \hfill
        \subcaptionbox{Модели профилей}{%
            \includegraphics[width=0.53\linewidth]{Рис_16b.jpg}}
        \hfill
    }
    \caption{Желоб Идзу-Бонин}\label{fig:16}
\end{figure}

Поперечные профили желоба Идзу-Бонин (ИБ) отображены как два сегмента ИБ (Рис. \ref{fig:16}). Результаты моделирования показали, что южный его сегмент более глубокий и имеет более выраженную V-образную форму с крутым уклоном на обоих склонах. Северная же часть желоба имеет асимметричную форму склона с заметной подводной террасой в западной части и выпрямленной формой со стороны океана. Кроме того, южный сегмент ИБ имеет на графике различимый хребет Бонин. Сравнение обоих сегментов ИБ со стороны океана указывает на схожую морфоструктуру, а также подводные горы в северной части ИБ, заметные на 3D-модели ИБ. Статистический анализ обоих сегментов показывает следующие результаты. Для северного сегмента, наиболее наглядный диапазон глубин; от -5000 до -6000 м где зарегистрировано 1037 записанных образцов (304, 407, 248 и 78), в то время как южный сегмент  показывает всего 931 образец (523, 304, 77 и 27). 

\pdfbookmark{Японский желоб}{jt} 
\underline{\textbf{Японский желоб}}. 
 ЯЖ простирается от Курильских о-вов до северной оконечности о. Идзу, достигая 8 046 м в глубине и 800 км в длине. ЯЖ соединяет Курило-Камчатский желоб на севере и желоб Идзу-Бонин на юге (Рис. \ref{fig:17}). ЯЖ образовался, при субдукции океанической Тихоокеанской плиты  под континентальную Охотскую плиту микроплиту. 

\begin{figure}[ht]
    \centerfloat{
        \hfill
        \subcaptionbox{Региональная батиметрия}{%
            \includegraphics[width=0.52\linewidth]{Рис_17a.jpg}}
        \hfill
        \subcaptionbox{Топографическая 3D модель}{%
            \includegraphics[width=0.42\linewidth]{Рис_17b.jpg}}
        \hfill
    }
    \caption{Японский желоб}\label{fig:17}
\end{figure}

Процесс субдукции вызывает кривизну изгиба опускающейся плиты, создавая условия для образования желоба и седиментации. Активные тектонические движения в зоне субдукции в районе ЯЖ продолжаются и отражаются в сейсмограммах, что является одной из основных причин цунами и землетрясений на севере Японии, в том числе мощного землетрясения Тохоку и вызванного им цунами, (11.03.2011). Скорость субдукции, связанная с ЯЖ, составляет около 7,9–9,2 см/год.

ЯЖ расположен в западной части Тихого океана, к востоку от Японии (Рис. \ref{fig:17}). Геологические и тектонические условия в районе ЯЖ сформированы субдукцией Тихоокеанской и Северо-Американской плит вдоль ЯЖ. Сравнительный анализ двух частей серии профилей поперечного сечения выявил более крутой склон в северной части желоба на континентальной части склона под фронтальным клином Северо-Американской плиты. Южная часть ЯЖ показывает более пологий склон на острове Хонсю. 

\begin{figure}[ht]
    \centerfloat{
        \hfill
        \subcaptionbox{Сегменты желоба}{%
            \includegraphics[width=0.45\linewidth]{Рис_18a.jpg}}
        \hfill
        \subcaptionbox{Модель профилей}{%
            \includegraphics[width=0.53\linewidth]{Рис_18b.jpg}}
        \hfill
    }
    \caption{Японский желоб}\label{fig:18}
\end{figure}

Расположение профилей, пересекающих два сегмента ЯЖ с юго-запада на северо-восток, показаны на карте (Рис. \ref{fig:18}). Было построено 12 профилей для южного сегмента и 14 профилей для северного. Геоморфологическое строение центральных профилей, и увеличение высоты в северном направлении видно в северной части. Результаты интерпретировались на основе гипотезы что геологические условия и сейсмичность района влияют на неровность рельефа желоба и общие геоморфологические формы при изучении батиметрических закономерностей и вариаций по профилям поперечного сечения ЯЖ, что показано на геологической карте с высоким уровнем сейсмичности, повторяющимися землетрясениями и вулканизмом вдоль желоба. 

Кроме того, значения, расположенные на глубинах от -4000 до -5000 м, чаще встречаются для южной части желоба: всего 134 точки, а северная часть желоба не превышает 100 значений: 91 пробы (287). Подводя итог, можно заключить, что южная часть желоба более мелкая чем северная. Она показывает более глубокие значения, и наиболее частое значение глубже, чем у северного сегмента. Анализ представленной модели геоида и гравитационные аномалий силы тяжести атмосферного воздуха в редукции Фая отражают изменения плотности Земли в области через ЯЖ и в его окрестностях. 

\pdfbookmark{Пюйсегюр желоб}{pst} 
\underline{\textbf{Желоб Пюйсегюр}}. Желоб Пюйсегюр простирается по дну южной части Тасманова моря, к югу от Новой Зеландии, более чем на 800 км к югу от ЮЗ точки побережья о. Южный (Рис. \ref{fig:19}). Географически, самая южная протяженность желоба зафиксирована на 400 километров к З от Оклендских о-вов. Тектонически он образован в результате субдукции Индо-Австралийской плиты под Тихоокеанскую. В востоку от желоба Пюйсегюр расположено северное продолжение хр. Маккуори. Район вокруг желоба Пюйсегюр характеризуется чрезвычайной сейсмической активностью и вулканизмом, что связано с Альпийским разломом, начинающимся на северном конце желоба.

\begin{figure}[ht]
    \centerfloat{
        \hfill
        \subcaptionbox{Топография желоба}{%
            \includegraphics[width=0.5\linewidth]{Рис_19a.jpg}}
        \hfill
        \subcaptionbox{Моделирование растров}{%
            \includegraphics[width=0.4\linewidth]{Рис_19b.jpg}}
        \hfill
    }
    \caption{Желоб Пюйсегюр}\label{fig:19}
\end{figure}

Желоб Пюйсегюр имеет асимметричную V-образную форму с крутым уклоном на восточном склоне и пологий склон на западном, соответственно (Рис. \ref{fig:19}). Прилегающий район имеет пологий склон на сторона океана для Хикуранги и крутой склон на восточном фланге для Пюйсегюра, окруженный сложным подводным рельефом вдоль дуги Маккуори. Пюйсегюр имеет бимодальный тип распределения данных с двумя четко выраженными вершинами.  Первый пик имеет три заметных диапазона c кластерами на диапазонах от -4200 до -4400 м (373 образца, частота данных 18\%), от -4000 до -4200 (182 образца, частота данных 9,6\%) и от -4400 м до -4600 м (216 образцов, частота данных 11,5\%).Вторая вершина покрывает диапазон с глубинами от -2200 до -3000 м с частотой данных ниже 7,5\%. Пик соответствует дуге Маккуори, граничащей со склоном желоба. 

Резкое уменьшение значений выборки с обеих сторон соответствует крутому склону желоба на западном его фланге от плато Хьорт. Абиссальная равнина простирается к востоку от Пюйсегюр, с небольшими выступами. Его морфоструктура особенно хорошо выражена вдоль дуги Маккуори, между 51,5° и 54° с четко выраженными гребнями, напоминающими вытянутые подводные горы на СЗ от Пюйсегюр, перемежающимися небольшими впадинами. В средней части желоба глубины резко увеличиваются от -5.400 до -6.200 м. Распределение данных имеет четкий бимодальный вид для Пюйсегюр с пиками в диапазонах от -4.200 до -4.400 (373 образца) и от -3.200 до -3.000 м (140 образцов). 

\begin{figure}[ht]
    \centerfloat{
        \hfill
        \subcaptionbox{Тектонико-геологические условия: Новая Зеландия}{%
            \includegraphics[width=0.42\linewidth]{Рис_20a.jpg}}
        \hfill
        \subcaptionbox{Моделирование профилей}{%
            \includegraphics[width=0.54\linewidth]{Рис_20b.jpg}}
        \hfill
    }
    \caption{Желоба Хикуранги, Пюйсегюр и Хьорт}\label{fig:20}
\end{figure}

\pdfbookmark{Хикуранги желоб}{ht} 
\underline{\textbf{Желоб Хикуранги}}. 
Желоб Хикуранги географически расположен у восточного побережья о. Северный Новой Зеландии, между южной оконечностью пролива Кука и возвышенностью Чатем. Тектонически, желоб представляет собой южное продолжение более глубокого желоба Кермадек, будучи расположенным в южном продолжении зоны субдукции Кермадек-Тонга. Желоб Хикуранги является зоной субдукции, где тяжелое океаническое плато Хикуранги погружается под континентальную кору Индо-Австралийской тектонической плиты. 

Активные геологические процессы (тектоническая эрозия, андерплейтинг) создают условия повышенной сейсмичности и вулканизма в районе желоба Хикуранги. Результаты проведенного моделирования желоба Хикуранги (ЖХ) выявили, что он имеет более корытообразную форму в поперечном сечении, в то время как желоба Пюйсегюр и Хьорт имеют классическую V-образную форму (Рис. \ref{fig:20}). Вариации морфоструктуры склонов были сопоставлены для желобов, что позволило выявить следующие результаты: ЖХ имеет корытообразную форму с плоским широким дном, также имеет более крутой уклон на склоне Северного острова. Форма гистограммы различается по всем трем желобам, отражающим их морфоструктуру. Распределение данных для ЖХ демонстрирует почти одинаковую картину для глубин -600 м до -2600 м с выборками менее 60 для каждого кластера (большинство данных при этом обнаружено на 30–40 выборках для желоба, но не более 4\%). Отчетливо видно возрастающее увеличение данных с глубинами ниже 500 м, то есть наклон в сторону о. Северного, частота от 5\% до 12\%. 

\FloatBarrier
\pdfbookmark{Заключение}{conclusion}                                  % Закладка pdf
В \underline{\textbf{заключении}} (Глава 5) приведено обсуждение основных результатов работы, которые заключаются в следующем.

Асимметрия глубоководных желобов, рассмотренная в настоящей работе, объясняется воздействием сложной совокупности тектонических, геологических и геодинамических факторов. Субдукция литосферных плит в мантию Земли вызывает образование океанических желобов, которые в конечном итоге формируются при заполнении пространства между плитами. Сложность и разнообразие форм рельефа в различных желобах возникает в результате влияния множества геологических процессов и факторов, влияющих на их форму, глубину, крутизны градиента и структуру. Из них тектонические условия, сейсмические и геологические условия региона являются решающими факторами. В этом можно убедиться, сравнив геофизические условия и морфологию дна желобов. Другие факторы, влияющие на морфологию желобов и их структуру, включают возраст и скорость конвергенции плит, промежуточные движения, ширину опускающейся плиты, толщину отложений на кристаллическом фундаменте, глубины, близость магматических вулканических областей и землетрясений, и географическое положение.

Глубоководные желоба являются важными сегментами океанического дна. Их морфоструктура отражает региональные тектонико-геологические процессы: движение, деформацию и изгиб плит, связанные с мантийной конвекцией. Распространение различных геологических факторов на прилегающих абиссальных равнинах океана способствует формированию геоморфологических особенностей структуры рельефа желобов. Комплексное изучение факторов является предпосылкой для понимания сложных процессов, формирующих желоба. На морфологию дна желобов влияют наносы потоками мутных течений, поступающих через сложную систему рельефа океанского дна. Тонкодисперсные осадки (глина, ил), откладываются из токов придонных течений, которые направляются континента по подводным каньонам и в конечном счете поступают в глубокие океанские желоба. Различные факторы влияют на формирование и морфоструктруты глубоководных желобов. Среди основных отметим движение тектонических плит, субдукцию низлежащих слоев плит в глубокую мантию, сейсмическая активность (подводные землетрясения). 

Широкий интерес к геологическому строению и тектонической эволюции океанических желобов, отмечен рядом публикаций. Представленная диссертация делает вклад в развитие этих исследований на примере 20 глубоководных тихоокеанских желобов. Моделирование, обработка, анализ и визуализация данных, представленные в данной работе внесли вклад в лучшее понимание геоморфологии желобов, а также обеспечили базу для дальнейших работ на основе выполненной серии тематических карт и визуализации морфологии 20 желобов. Факторы, определяющие формирование их морфоструктур, были проанализированы. Отмечено, что тектонический фактор (скорость движения литосферных плит, плотность и возраст литосферы в зоне субдукции) является решающим

%% Согласно ГОСТ Р 7.0.11-2011:
%% 5.3.3 В заключении диссертации излагают итоги выполненного исследования, рекомендации, перспективы дальнейшей разработки темы.
%% 9.2.3 В заключении автореферата диссертации излагают итоги данного исследования, рекомендации и перспективы дальнейшей разработки темы.
%\begin{enumerate}
  %\item На основе анализа \ldots
  %\item Численные исследования показали, что \ldots
  %\item Математическое моделирование показало \ldots
  %\item Для выполнения поставленных задач был создан \ldots
%\end{enumerate}

Представленная диссертация предлагает обзор морфоструктур 20 желобов Тихого океана, выполненный с использованием передовых методов анализа данных, моделирования и картографирования. Исследование представляет собой междисциплинарный подход, сочетающий картографическое моделирование и тектонико-геологический региональный анализ. Результаты позволили представить вариативность желобов на основе их геоморфологических свойств, изменения уклонов, анализ градиентов крутизны и характер батиметрии. Анализ данных в геотектонике и геологии с использованием статистических и скриптовых картографических методов является важным подходом, направленным на повышение точности обработки больших массивов данных. Методы, представленные в диссертации, направлены на решение проблемы анализа датасетов для вопросов моделирования в морской геоморфологии, что связано с тем, что существующие решения в основном применяют традиционные методы ГИС. Диссертация продемонстрировала использование методов скриптовой картографии применительно к задачам геотектоники и геологии для  анализа подводных структур рельефа дна на примере 20 желобов Тихого океана. 


\pdfbookmark{Литература}{bibliography}                                % Закладка pdf

\ifdefmacro{\microtypesetup}{\microtypesetup{protrusion=false}}{} % не рекомендуется применять пакет микротипографики к автоматически генерируемому списку литературы
\urlstyle{rm}                               % ссылки URL обычным шрифтом
\ifnumequal{\value{bibliosel}}{0}{% Встроенная реализация с загрузкой файла через движок bibtex8
    \renewcommand{\bibname}{\large \bibtitleauthor}
    \nocite{*}
    \insertbiblioauthor           % Подключаем Bib-базы
    %\insertbiblioexternal   % !!! bibtex не умеет работать с несколькими библиографиями !!!
}{% Реализация пакетом biblatex через движок biber
    % Цитирования.
    %  * Порядок перечисления определяет порядок в библиографии (только внутри подраздела, если `\insertbiblioauthorgrouped`).
    %  * Если не соблюдать порядок "как для \printbibliography", нумерация в `\insertbiblioauthor` будет кривой.
    %  * Если цитировать каждый источник отдельной командой --- найти некоторые ошибки будет проще.
    %
    %% authorvak   
    	
    		\nocite{Lemenkova202105}
    		\nocite{Lemenkova202104}
		\nocite{Lemenkova202032}	
		\nocite{Lemenkova202027}
		\nocite{Lemenkova202018}
		\nocite{Lemenkova202017}
		\nocite{Lemenkova202016}
		\nocite{Lemenkova202015}
		\nocite{Lemenkova202012}
		\nocite{Lemenkova202005}
		\nocite{Lemenkova202003}
		\nocite{Lemenkova201995}
		\nocite{Lemenkova201994}
		\nocite{Lemenkova201991}	
		\nocite{Lemenkova201904}
		\nocite{Lemenkova201901}						
		\nocite{Леменкова2015a}
		\nocite{Леменкова2015b}
		\nocite{Суетова2005}

    \ifnumgreater{\value{usefootcite}}{0}{
        \begin{refcontext}[labelprefix={}]
            \ifnum \value{bibgrouped}>0
                \insertbiblioauthorgrouped    % Вывод всех работ автора, сгруппированных по источникам
            \else
                \insertbiblioauthor      % Вывод всех работ автора
            \fi
        \end{refcontext}
    }{
        \ifnum \totvalue{citeexternal}>0
            \begin{refcontext}[labelprefix=A]
                \ifnum \value{bibgrouped}>0
                    \insertbiblioauthorgrouped    % Вывод всех работ автора, сгруппированных по источникам
                \else
                    \insertbiblioauthor      % Вывод всех работ автора
                \fi
            \end{refcontext}
        \else
            \ifnum \value{bibgrouped}>0
                \insertbiblioauthorgrouped    % Вывод всех работ автора, сгруппированных по источникам
            \else
                \insertbiblioauthor      % Вывод всех работ автора
            \fi
        \fi
        %  \insertbiblioauthorimportant  % Вывод наиболее значимых работ автора (определяется в файле characteristic во второй section)
        \begin{refcontext}[labelprefix={}]
            \insertbiblioexternal            % Вывод списка литературы, на которую ссылались в тексте автореферата
        \end{refcontext}
        % Невидимый библиографический список для подсчёта количества внешних публикаций
        % Используется, чтобы убрать приставку "А" у работ автора, если в автореферате нет
        % цитирований внешних источников.
        \printbibliography[heading=nobibheading, section=0, env=countexternal, keyword=biblioexternal, resetnumbers=true]%
    }
}
\ifdefmacro{\microtypesetup}{\microtypesetup{protrusion=true}}{}
\urlstyle{tt}                               % возвращаем установки шрифта ссылок URL
